\documentclass[11pt]{article}
\usepackage{mathblatt}
\begin{document}
\blattkopf*{Basisaufgaben}{BAS-Z7}{Basisaufgaben · Zettel · BAS-Z7}
\einheitenkopf[][Zettel 7]{Basisaufgaben · Zettel 7}
\zweigzeile{ohne Taschenrechner · je Aufgabe 1 Punkt · Lösungen auf der Rückseite}
% \small: zehn Aufgaben mit bis zu zwei Koordinatensystemen auf einer Seite (Render 28.09.)
\small

% 1: Bruchteil einer Fläche bestimmen – brueche-dezimalzahlen-basis-k1-v6 (Original 2026-FOR-B1b)
\begin{aufgabe}{Bei welcher Figur ist genau $\frac{2}{5}$ grau? \hfill (P26F) \\ \kreuz{P} \kreuz{Q} \kreuz{R} \kreuz{S}}

P \kreissektor[0.8]{144}{} \quad Q \bruchkreis[0.8]{2}{6} \quad R \bruchrechteck[2.5]{5}{10} \quad S \bruchrechteck[2.5]{2}{4}
\end{aufgabe}

% 2: Zahlen in verschiedenen Darstellungen vergleichen – brueche-dezimalzahlen-basis-k4-v6 (Original 2023-OS-B1f)
\begin{aufgabe}{Welcher Wert ist der kleinste: $0{,}5$; $0{,}5^2$; $5\,\%$; $\frac{1}{5}$? \hfill (P23) \\ \leerfeld}
\end{aufgabe}

% 3: Zehnerpotenzschreibweise umwandeln – potenzen-wurzeln-basis-k3-v3 (Original 2025-OS-B1d)
\begin{aufgabe}{$0{,}0045 = 4{,}5 \cdot 10^{\square}$ – welche Hochzahl? \hfill (P25) \\ \leerfeld}
\end{aufgabe}

% 4: Prozent und Anteil umwandeln – prozentrechnung-basis-k2-v2 (Original 2022-OS-B1f)
\begin{aufgabe}{Jedes zweite Kind hat ein Haustier – wie viel Prozent? Kreuze an. \hfill (P22) \\ \kreuz{$2\,\%$} \kreuz{$20\,\%$} \kreuz{$50\,\%$} \kreuz{$200\,\%$}}
\end{aufgabe}

% 5: Uhrzeit aus Startzeit und Dauer berechnen – einheiten-basis-k1-v1 (Original 2021-OS-B1a)
\begin{aufgabe}{Ein Zug fährt um 8:45 Uhr ab und ist $2$ h $30$ min unterwegs. Wann kommt er an? \hfill (P21) \\ \leerfeld Uhr}
\end{aufgabe}

% 6: Lineare Gleichung lösen – lineare-gleichungen-basis-k1-v4 (Original 2024-OS-B1d)
\begin{aufgabe}{Löse die Gleichung $2 \cdot (x - 5) + 7 = 7$. \hfill (P24) \\ x = \leerfeld}
\end{aufgabe}

% 7: Lösung durch Einsetzen prüfen – quadratische-gleichungen-basis-k1-v3 (Original 2025-OS-B1h)
\begin{aufgabe}{Welcher Wert erfüllt $x \cdot (x - 3) = 10$? \hfill (P25) \\ \kreuz{$x = 2$} \kreuz{$x = 5$} \kreuz{$x = -5$} \kreuz{$x = 10$}}
\end{aufgabe}

% 8: Symmetrieachsen bestimmen – symmetrie-abbildungen-basis-k3-v2 (Original 2022-OS-B1i)
\begin{aufgabe}{\begin{minipage}[t]{0.6\linewidth}Eine Fliese hat die Form einer Raute (kein Quadrat). Wie viele Symmetrieachsen hat die Figur? Kreuze an. \hfill (P22) \\ \kreuz{0} \kreuz{1} \kreuz{2} \kreuz{3} \kreuz{4} \kreuz{5}\end{minipage}\hfill\begin{minipage}[t]{0.37\linewidth}\vspace{0pt}
\raute{2.8}{1.6}
\end{minipage}}
\end{aufgabe}

% 9: Volumen Würfel berechnen – koerper-basis-k3-v1 (Original 2026-FOR-B1f)
\begin{aufgabe}{Ein Würfel hat die Kantenlänge $4$ cm. Gib sein Volumen an. \hfill (P26F) \\ \leerfeld[cm³]}
\end{aufgabe}

% 10: Winkelfunktion Seitenverhältnis angeben – trigonometrie-basis-k2-v4 (Original 2025-OS-B1g)
\begin{aufgabe}{\begin{minipage}[t]{0.6\linewidth}Rechter Winkel bei $B$. Trage den Bruch ein. \hfill (P25) \\ $\mathrm{cos}\,\gamma =$ \leerfeld\end{minipage}\hfill\begin{minipage}[t]{0.37\linewidth}\vspace{0pt}
\dreieck{(0,0)}{(3,0)}{(3,2.2)}{g}{h}{i}{}{}{\gamma}
\end{minipage}}
\end{aufgabe}

\begleitteil
\erg{1}{P, denn $\frac{144}{360} = \frac{2}{5}$}
\erg{2}{$5\,\% = 0{,}05$; die anderen sind $0{,}5$, $0{,}25$ und $0{,}2$}
\erg{3}{$-3$ – das Komma wandert drei Stellen nach rechts}
\erg{4}{$50\,\%$ ($\frac{1}{2} = \frac{50}{100}$)}
\erg{5}{8:45 Uhr $+ 2$ h $=$ 10:45 Uhr, $+ 30$ min $=$ 11:15 Uhr}
\erg{6}{$2 \cdot (x - 5) = 0$, also $x = 5$}
\erg{7}{$x = 5$, denn $5 \cdot (5 - 3) = 5 \cdot 2 = 10$}
\erg{8}{2}
\erg{9}{$4 \cdot 4 \cdot 4 = 64$ cm³ (nicht $4 \cdot 4 = 16$)}
\erg{10}{$\mathrm{cos}\,\gamma = \frac{g}{h}$}
\end{document}
